%% toc: P3 40 40.2
\section{Proofs: random graphs and extension statements}
\label{sec:annex-p03-w03-p03-065-random-graphs}

Throughout, \(G(n,p)\) is the random graph of
\cref{def:p03-w03-p03-065-gnp}, \(\ln\) is the natural logarithm, and two
elementary inequalities are used repeatedly.  First, \(1-t\leq e^{-t}\) for
every real \(t\), since \(e^{-t}\) is convex and \(1-t\) is its tangent line
at \(t=0\).  Second, \(\ln(1-t)\geq-t-t^2\) for \(0\leq t\leq1/2\): the
function \(h(t)=\ln(1-t)+t+t^2\) has \(h(0)=0\) and
\[
  h'(t)=-\frac1{1-t}+1+2t=\frac{t(1-2t)}{1-t}\geq0
  \qquad(0\leq t\leq1/2).
\]

\begin{annexitem}{annex:p03-w03-p03-065-moments}{thm:p03-w03-p03-065-moments}{First and second moment bounds}
\begin{proof}
Since \(X\) takes values in \(\{0,1,2,\ldots\}\), the indicator of
\(\{X>0\}\) is at most \(X\) pointwise.  Taking expectations gives
\(\Pr[X>0]\leq\mathbb E[X]\).

Let \(m=\mathbb E[X]>0\).  If \(X=0\) then \(|X-m|\geq m\), so
\(\mathbf 1_{\{X=0\}}\leq(X-m)^2/m^2\) pointwise.  Taking expectations gives
\(\Pr[X=0]\leq\operatorname{Var}[X]/m^2\).

For \(X=\sum_{i\in I}\mathbf 1_{B_i}\), bilinearity of covariance gives
\[
  \operatorname{Var}[X]=\sum_{i\in I}\bigl(\Pr[B_i]-\Pr[B_i]^2\bigr)
  +\sum_{i\neq j}\bigl(\Pr[B_i\wedge B_j]-\Pr[B_i]\Pr[B_j]\bigr).
\]
The first sum is at most \(\sum_i\Pr[B_i]=\mathbb E[X]\).  In the second
sum, every ordered pair outside \(D\) has \(B_i\) and \(B_j\) independent, so
its term is \(0\).  This leaves the sum over \(D\).
\end{proof}
\end{annexitem}

\begin{annexitem}{annex:p03-w03-p03-065-clique}{thm:p03-w03-p03-065-clique}{Threshold for a complete subgraph}
\begin{proof}
Write \(m=\binom k2\).  For each \(k\)-element set \(S\subseteq[n]\) let
\(A_S\) be the event that all \(m\) pairs inside \(S\) are edges, so
\(\Pr[A_S]=p^m\).  Let \(X=\sum_S\mathbf 1_{A_S}\), the number of \(k\)-sets
spanning a complete graph.  Then \(G(n,p)\entails A_k\) exactly when
\(X>0\), and
\[
  \mathbb E[X]=\binom nk p^m .
\]
By \cref{def:p03-w03-p03-065-asymptotic-order} with
\(\gamma=(k-1)/2\), \(p\ll n^{-2/(k-1)}\) holds exactly when
\(np^{(k-1)/2}\to0\), and \(p\gg n^{-2/(k-1)}\) holds exactly when
\(p>0\) eventually and \(np^{(k-1)/2}\to\infty\).

\emph{Below the threshold.}  Since \(\binom nk\leq n^k\) and
\(p^m=(p^{(k-1)/2})^k\),
\(\mathbb E[X]\leq(np^{(k-1)/2})^k\to0\).  By
\cref{thm:p03-w03-p03-065-moments}, \(\Pr[X>0]\to0\).

\emph{Above the threshold.}  Let \(n\geq k\) and \(p>0\).  Since
\(\binom nk\geq(n/k)^k\),
\(\mathbb E[X]\geq k^{-k}(np^{(k-1)/2})^k\to\infty\), and in particular
\(\mathbb E[X]>0\).  Two events \(A_S,A_T\) with \(|S\cap T|\leq1\) concern
disjoint sets of pairs, so they are independent.  Apply
\cref{thm:p03-w03-p03-065-moments} with \(D\) the set of ordered pairs
\((S,T)\) with \(2\leq|S\cap T|\leq k-1\).  For such a pair with
\(|S\cap T|=i\), the pairs inside \(S\) or inside \(T\) number
\(2m-\binom i2\), so \(\Pr[A_S\wedge A_T]=p^{2m-\binom i2}\).  The number of
ordered pairs \((S,T)\) with \(|S\cap T|=i\) is
\(\binom nk\binom ki\binom{n-k}{k-i}\leq\binom nk\binom ki n^{k-i}\).  Hence
\[
  \frac{\operatorname{Var}[X]}{\mathbb E[X]^2}
  \leq\frac1{\mathbb E[X]}
  +\sum_{i=2}^{k-1}\frac{\binom nk\binom ki n^{k-i}p^{2m-\binom i2}}
                        {\binom nk^2p^{2m}}
  \leq\frac1{\mathbb E[X]}
  +\sum_{i=2}^{k-1}\binom ki\,k^k\,\frac1{n^ip^{\binom i2}} .
\]
For \(2\leq i\leq k-1\),
\(n^ip^{\binom i2}=(np^{(i-1)/2})^i\geq(np^{(k-1)/2})^i\), because
\(0<p\leq1\) and \((i-1)/2\leq(k-1)/2\).  The right-hand side tends to
\(\infty\), so every term of the finite sum tends to \(0\).  By
\cref{thm:p03-w03-p03-065-moments}, \(\Pr[X=0]\to0\).  For \(k=2\) the sum
is empty and the same conclusion holds.
\end{proof}
\end{annexitem}

\begin{annexitem}{annex:p03-w03-p03-065-isolated}{thm:p03-w03-p03-065-isolated}{Threshold for having no isolated vertex}
\begin{proof}
Let \(X\) be the number of isolated vertices, so that \(I\) holds exactly
when \(X=0\).  A vertex is isolated when none of its \(n-1\) pairs is an
edge, so \(\mathbb E[X]=n(1-p)^{n-1}\).

\emph{Part 1.}  Using \(1-p\leq e^{-p}\) and \(p(n-1)=pn-p\geq
(1+\varepsilon)\ln n-1\),
\[
  \mathbb E[X]\leq n\,e^{-p(n-1)}\leq e\,n\cdot n^{-(1+\varepsilon)}
  =e\,n^{-\varepsilon}\to0 ,
\]
so \(\Pr[X>0]\to0\) by \cref{thm:p03-w03-p03-065-moments}.

\emph{Part 2.}  Here \(p\leq(1-\varepsilon)\ln n/n\to0\), so \(p\leq1/2\)
for all large \(n\).  Using \(\ln(1-p)\geq-p-p^2\),
\[
  \mathbb E[X]\geq n\,e^{-(p+p^2)(n-1)}\geq n\,e^{-pn}e^{-p^2n}
  \geq n^{\varepsilon}e^{-p^2n},
\]
and \(p^2n\leq(\ln n)^2/n\to0\).  So \(\mathbb E[X]\to\infty\).  Two
distinct vertices \(u,w\) are both isolated when none of the \(2n-3\) pairs
containing \(u\) or \(w\) is an edge.  Hence
\(\mathbb E[X(X-1)]=n(n-1)(1-p)^{2n-3}\) and
\[
  \frac{\operatorname{Var}[X]}{\mathbb E[X]^2}
  =\frac{\mathbb E[X(X-1)]+\mathbb E[X]-\mathbb E[X]^2}{\mathbb E[X]^2}
  =\frac1{\mathbb E[X]}+\frac{n-1}{n(1-p)}-1 ,
\]
which is at most \(1/\mathbb E[X]+p/(1-p)\), and this tends to \(0\):
\[
  \frac{\operatorname{Var}[X]}{\mathbb E[X]^2}
  \leq\frac1{\mathbb E[X]}+\frac p{1-p}\to0 .
\]
By \cref{thm:p03-w03-p03-065-moments}, \(\Pr[X=0]\to0\).

\emph{Threshold.}  If \(p\gg n^{-1}\ln n\), then \(pn\geq2\ln n\) for all
large \(n\), and part~1 applies with \(\varepsilon=1\).  If
\(p\ll n^{-1}\ln n\), then \(pn\leq\frac12\ln n\) for all large \(n\), and
part~2 applies with \(\varepsilon=1/2\).
\end{proof}
\end{annexitem}

\begin{annexitem}{annex:p03-w03-p03-065-connectivity}{thm:p03-w03-p03-065-connectivity}{Connectivity in the critical window}
\begin{proof}
Write \(L=np=\ln n+c\) and \(\lambda=e^{-c}\).  Take \(n\) so large that
\(1\leq L\leq2\ln n\) and \(p\leq1/2\).

\emph{Factorial moments of \(X\).}  For \(k\geq1\) let \((X)_k=X(X-1)\cdots
(X-k+1)\), the number of \(k\)-tuples of distinct isolated vertices.  A tuple
of \(k\) distinct vertices consists of isolated vertices exactly when none of
the \(M_k=k(n-k)+\binom k2\) pairs meeting the tuple is an edge, so
\(\mathbb E[(X)_k]=n(n-1)\cdots(n-k+1)\,(1-p)^{M_k}\).  Fix \(k\).  Since
\(M_k=kn-(k^2+k)/2\) and \(-p-p^2\leq\ln(1-p)\leq-p\),
\[
  -pkn-p^2kn\leq M_k\ln(1-p)\leq-pkn+p(k^2+k)/2 .
\]
Here \(pkn=k\ln n+kc\), while \(p^2kn=kL^2/n\to0\) and
\(p(k^2+k)/2\to0\).  So \(n^k(1-p)^{M_k}\to e^{-kc}\), and since
\(n(n-1)\cdots(n-k+1)/n^k\to1\), \(\mathbb E[(X)_k]\to\lambda^k\).  Put
\((X)_0=1\).

\emph{Alternating bounds.}  For integers \(N\geq1\) and \(j\geq0\),
\[
  \sum_{i=0}^{j}(-1)^i\binom Ni=(-1)^j\binom{N-1}j .
\]
For \(j=0\) both sides are \(1\).  If it holds for \(j\), adding
\((-1)^{j+1}\binom N{j+1}\) gives
\((-1)^{j+1}\bigl(\binom N{j+1}-\binom{N-1}j\bigr)=(-1)^{j+1}\binom{N-1}{j+1}\)
by Pascal's rule.  For \(N=0\) the sum is \(1\).  Hence the partial sum is at
least \(\mathbf 1_{\{N=0\}}\) when \(j\) is even and at most
\(\mathbf 1_{\{N=0\}}\) when \(j\) is odd.  Put \(N=X\) and take
expectations, using \(\mathbb E\binom Xi=\mathbb E[(X)_i]/i!\): for every
\(m\geq0\),
\[
  \sum_{i=0}^{2m+1}\frac{(-1)^i\,\mathbb E[(X)_i]}{i!}\leq\Pr[X=0]
  \leq\sum_{i=0}^{2m}\frac{(-1)^i\,\mathbb E[(X)_i]}{i!}.
\]
Letting \(n\to\infty\) with \(m\) fixed, the limit inferior of \(\Pr[X=0]\)
is at least \(\sum_{i\leq2m+1}(-\lambda)^i/i!\) and the limit superior is at
most \(\sum_{i\leq2m}(-\lambda)^i/i!\).  Both partial sums tend to
\(e^{-\lambda}\) as \(m\to\infty\), so \(\Pr[X=0]\to e^{-\lambda}\).

\emph{Trees.}  A connected graph on a \(k\)-element set \(S\) has a spanning
tree: a connected spanning subgraph with the fewest edges has no cycle,
because deleting an edge of a cycle leaves it connected.  The trees on \(S\)
number at most \(k^{k-1}\): fix \(s_0\in S\) and send a tree to the map
\(S\setminus\{s_0\}\to S\) taking each vertex to its neighbour on the tree
path to \(s_0\); the tree's edges are the pairs \(\{s,\varphi(s)\}\), so the
tree is recovered from the map.  The same pairs show that a tree on \(S\) has
\(k-1\) edges.  Hence \(\Pr[G(n,p)|_S\text{ connected}]\leq k^{k-1}p^{k-1}\)
by the union bound over spanning trees.

\emph{No middle components.}  Let \(Y\) be the number of connected
components with between \(2\) and \(n/2\) vertices.  A set \(S\) of size
\(k\) is the vertex set of a component when \(G(n,p)|_S\) is connected and
none of the \(k(n-k)\) pairs between \(S\) and its complement is an edge;
these concern disjoint sets of pairs, so are independent.  Using
\(\binom nk\leq n^k/k!\leq(en/k)^k\) (because \(e^k\geq k^k/k!\)) and
\(1-p\leq e^{-p}\),
\[
  \Pr[Y>0]\leq\sum_{k=2}^{\lfloor n/2\rfloor}\binom nk k^{k-1}p^{k-1}
  (1-p)^{k(n-k)}
  \leq\sum_{k=2}^{\lfloor n/2\rfloor}\frac1{kp}
  \bigl(e\,np\,e^{-p(n-k)}\bigr)^k .
\]
For \(2\leq k\leq\sqrt n\): \(e^{-p(n-k)}=e^{-L}e^{pk}=(\lambda/n)e^{pk}\)
and \(pk\leq L/\sqrt n\leq1\) for large \(n\), so the base is at most
\(DL/n\) with \(D=e^2\lambda\); and \(1/(kp)\leq n/(2L)\).  For large \(n\),
\(DL/n\leq1/2\), and these terms sum to at most
\[
  \frac n{2L}\sum_{k\geq2}\Bigl(\frac{DL}n\Bigr)^k
  \leq\frac n{2L}\cdot2\Bigl(\frac{DL}n\Bigr)^2=\frac{D^2L}n\to0 .
\]
For \(\sqrt n<k\leq n/2\): \(n-k\geq n/2\), so
\(e^{-p(n-k)}\leq e^{-L/2}=\lambda^{1/2}n^{-1/2}\), and the base is at most
\(e\lambda^{1/2}Ln^{-1/2}\leq n^{-1/3}\) for large \(n\); also
\(1/(kp)\leq n/L\leq n\).  Each of these at most \(n\) terms is at most
\(n\cdot n^{-\sqrt n/3}\), so they sum to at most \(n^2n^{-\sqrt n/3}\to0\).
Hence \(\Pr[Y>0]\to0\).

\emph{Conclusion.}  For \(n\geq2\), a connected graph has no isolated
vertex, so \(\Pr[\text{connected}]\leq\Pr[X=0]\).  If \(X=0\) and the graph
is not connected, it has at least two components, and the smallest has at
most \(n/2\) vertices and, having no isolated vertex, at least \(2\); so
\(Y>0\).  Hence \(\Pr[\text{connected}]\geq\Pr[X=0]-\Pr[Y>0]\).  Both bounds
tend to \(e^{-\lambda}=e^{-e^{-c}}\).
\end{proof}
\end{annexitem}

\begin{annexitem}{annex:p03-w03-p03-065-common-neighbour}{thm:p03-w03-p03-065-common-neighbour}{Threshold for common neighbours}
\begin{proof}
For an unordered pair \(P=\{a,b\}\) of distinct vertices, let \(A_P\) be the
event that \(a\) and \(b\) have no common neighbour.  For distinct
\(z\notin P\) the events \(z\sim a\wedge z\sim b\) concern disjoint pairs, so
they are independent, and each has probability \(p^2\).  Hence
\(\Pr[A_P]=(1-p^2)^{n-2}\).

\emph{Part 1.}  By the union bound over at most \(n^2/2\) pairs,
\[
  \Pr[\neg C]\leq\frac{n^2}2\,e^{-p^2(n-2)}
  =\frac{n^2}2\,e^{-p^2n}e^{2p^2}
  \leq\frac{e^2}2\,n^2n^{-(2+\varepsilon)}\to0 .
\]

\emph{Part 2.}  Now \(p^2\leq(2-\varepsilon)\ln n/n\to0\); take \(n\) so
large that \(p^2\leq1/2\).  Let \(Y=\sum_P\mathbf 1_{A_P}\), so that \(C\)
holds exactly when \(Y=0\).  Using \(\binom n2\geq n^2/4\) and
\(\ln(1-p^2)\geq-p^2-p^4\),
\[
  \mathbb E[Y]=\binom n2(1-p^2)^{n-2}
  \geq\frac{n^2}4\,e^{-p^2n-p^4n}
  \geq\frac14\,n^{\varepsilon}e^{-p^4n}\to\infty ,
\]
since \(p^4n\leq(2\ln n)^2/n\to0\).  Apply
\cref{thm:p03-w03-p03-065-moments} with \(D\) the set of all ordered pairs
\((P,Q)\) with \(P\neq Q\).  These split into two kinds.

\emph{Disjoint pairs} \(P=\{a,b\}\), \(Q=\{c,d\}\).  Let
\(W=[n]\setminus\{a,b,c,d\}\).  The event \(A_P\wedge A_Q\) implies that, for
every \(z\in W\), neither \(z\sim a\wedge z\sim b\) nor
\(z\sim c\wedge z\sim d\).  For different \(z\) these conditions concern
disjoint sets of pairs, and for one \(z\) the two conditions concern
disjoint pairs.  So
\[
  \Pr[A_P\wedge A_Q]\leq(1-p^2)^{2(n-4)}=(1-p^2)^{-4}\Pr[A_P]\Pr[A_Q].
\]
Summing over these ordered pairs,
\[
  \sum\bigl(\Pr[A_P\wedge A_Q]-\Pr[A_P]\Pr[A_Q]\bigr)
  \leq\bigl((1-p^2)^{-4}-1\bigr)\mathbb E[Y]^2 .
\]

\emph{Pairs sharing one vertex}, \(P=\{a,b\}\), \(Q=\{a,c\}\) with
\(b\neq c\).  Let \(W=[n]\setminus\{a,b,c\}\).  For \(z\in W\), the event
that \(z\) is a common neighbour of \(a,b\) or of \(a,c\) is
\(z\sim a\wedge(z\sim b\vee z\sim c)\), of probability
\(p(1-(1-p)^2)=2p^2-p^3\).  These events are independent for different
\(z\), so
\(\Pr[A_P\wedge A_Q]\leq q^{n-3}\) with \(q=1-2p^2+p^3\).  There are at most
\(n^3\) such ordered pairs.  Since \(\mathbb E[Y]^2\geq\frac{n^4}{16}
(1-p^2)^{2n-4}\),
\[
  \frac{n^3q^{n-3}}{\mathbb E[Y]^2}
  \leq\frac{16}n\Bigl(\frac q{(1-p^2)^2}\Bigr)^{n-3}(1-p^2)^{-2}.
\]
Here \((1-p^2)^2=1-2p^2+p^4\geq1/4\), so
\(q/(1-p^2)^2=1+(p^3-p^4)/(1-p^2)^2\leq1+4p^3\leq e^{4p^3}\), and
\(p^3n=p\cdot p^2n\leq2\ln n\,\sqrt{2\ln n/n}\to0\).  So the displayed bound
is at most \((16/n)\,e^{4p^3n}\cdot4\to0\).

Combining the two kinds,
\[
  \Pr[Y=0]\leq\frac{\operatorname{Var}[Y]}{\mathbb E[Y]^2}
  \leq\frac1{\mathbb E[Y]}+\bigl((1-p^2)^{-4}-1\bigr)
  +\frac{n^3q^{n-3}}{\mathbb E[Y]^2}\to0 ,
\]
since \(p\to0\).

\emph{Threshold.}  If \(p\gg n^{-1/2}\ln^{1/2}n\), then \(p^2n/\ln n\to
\infty\), so \(p^2n\geq3\ln n\) eventually and part~1 applies with
\(\varepsilon=1\).  If \(p\ll n^{-1/2}\ln^{1/2}n\), then
\(p^2n\leq\ln n\) eventually and part~2 applies with \(\varepsilon=1\).  At
\(p=n^{-1/2}\), \(p^2n=1\leq\ln n\) for \(n\geq3\), so part~2 applies again.
\end{proof}
\end{annexitem}

\begin{annexitem}{annex:p03-w03-p03-065-harris}{thm:p03-w03-p03-065-harris}{Harris inequality}
\begin{proof}
Induction on \(m\).  For \(m=0\) the functions are constants and equality
holds.  Let \(m\geq1\), write \(\omega=(\omega',\omega_m)\) with
\(\omega'=(\omega_1,\ldots,\omega_{m-1})\), and let
\(\pi_t=\Pr[\omega_m=t]\) for \(t\in\{0,1\}\).

\emph{One coordinate.}  Fix \(\omega'\) and put \(\varphi(t)=f(\omega',t)\),
\(\psi(t)=g(\omega',t)\).  Then \(\varphi\) is nondecreasing and \(\psi\)
nonincreasing on \(\{0,1\}\), so
\((\varphi(s)-\varphi(t))(\psi(s)-\psi(t))\leq0\) for all \(s,t\).
Multiplying by \(\pi_s\pi_t\) and summing over \(s,t\in\{0,1\}\) gives
\[
  2\Bigl(\sum_t\pi_t\varphi(t)\psi(t)
  -\sum_t\pi_t\varphi(t)\sum_t\pi_t\psi(t)\Bigr)\leq0 .
\]
Define \(f_1(\omega')=\sum_t\pi_tf(\omega',t)\) and
\(g_1(\omega')=\sum_t\pi_tg(\omega',t)\).  Since \(\omega_m\) is independent
of \(\omega'\), the display says
\(\mathbb E[fg\mid\omega']\leq f_1(\omega')g_1(\omega')\).

\emph{Induction.}  The function \(f_1\) is a combination, with nonnegative
weights, of functions nondecreasing in each coordinate of \(\omega'\), so it
is nondecreasing; likewise \(g_1\) is nonincreasing.  By the induction
hypothesis applied to \(\omega'\),
\[
  \mathbb E[fg]=\mathbb E\bigl[\mathbb E[fg\mid\omega']\bigr]
  \leq\mathbb E[f_1g_1]\leq\mathbb E[f_1]\,\mathbb E[g_1]
  =\mathbb E[f]\,\mathbb E[g].
\]
For events, \(B\) is increasing exactly when \(\mathbf 1_B\) is
nondecreasing in each coordinate, and \(D\) is decreasing exactly when
\(\mathbf 1_D\) is nonincreasing; apply the inequality to these indicators.
\end{proof}
\end{annexitem}

\begin{annexitem}{annex:p03-w03-p03-065-janson}{thm:p03-w03-p03-065-janson}{Janson inequality}
\begin{proof}
If some \(E_i\) is empty then \(B_i\) is certain, the left-hand side is
\(0\), and both bounds hold.  Assume every \(E_i\) is nonempty and enumerate
\(\Omega=\{1,\ldots,N\}\).  The configuration with every coordinate \(0\)
has positive probability, because \(\Pr[\omega_k=0]>0\) for each \(k\), and
it lies in every \(\neg B_j\).  So every intersection of events \(\neg B_j\)
has positive probability, and the conditional probabilities below are
defined.

\emph{The first bound.}  Let \(r_i=\Pr[B_i\mid\bigwedge_{j<i}\neg B_j]\).
Then
\[
  \Pr\Bigl[\bigwedge_{i=1}^N\neg B_i\Bigr]=\prod_{i=1}^N(1-r_i).
\]
Fix \(i\).  Let \(D_1\) be the intersection of the \(\neg B_j\) with
\(j<i\) and \(j\approx i\), and \(D_2\) the intersection of those with \(j<i\)
and \(j\not\approx i\) (an empty intersection is the sure event).  Then
\[
  r_i=\frac{\Pr[B_i\wedge D_1\wedge D_2]}{\Pr[D_1\wedge D_2]}
  \geq\frac{\Pr[B_i\wedge D_1\wedge D_2]}{\Pr[D_2]}
  =\Pr[B_i\wedge D_1\mid D_2],
\]
and
\[
  \Pr[B_i\wedge D_1\mid D_2]
  \geq\Pr[B_i\mid D_2]-\sum_{j<i,\;j\approx i}\Pr[B_i\wedge B_j\mid D_2],
\]
the last step because \(B_i\wedge\neg D_1\) is contained in the union of the
events \(B_i\wedge B_j\) with \(j<i\), \(j\approx i\).  The event \(B_i\)
depends only on the coordinates in \(E_i\), and \(D_2\) only on coordinates
in sets \(E_j\) disjoint from \(E_i\), so they are independent and
\(\Pr[B_i\mid D_2]=\Pr[B_i]\).  The event \(B_i\wedge B_j\) is increasing and
\(D_2\) is decreasing, so by \cref{thm:p03-w03-p03-065-harris}
\(\Pr[B_i\wedge B_j\wedge D_2]\leq\Pr[B_i\wedge B_j]\Pr[D_2]\), that is,
\(\Pr[B_i\wedge B_j\mid D_2]\leq\Pr[B_i\wedge B_j]\).  Hence
\[
  r_i\geq\Pr[B_i]-\sum_{j<i,\;j\approx i}\Pr[B_i\wedge B_j].
\]
Using \(1-r_i\leq e^{-r_i}\),
\[
  \Pr\Bigl[\bigwedge_i\neg B_i\Bigr]\leq\exp\Bigl(-\sum_ir_i\Bigr)
  \leq\exp\Bigl(-\mu+\sum_i\sum_{j<i,\;j\approx i}\Pr[B_i\wedge B_j]\Bigr)
  =e^{-\mu+\Delta/2},
\]
since each unordered pair \(\{i,j\}\) with \(i\approx j\) occurs once in the
double sum and twice in \(\Delta\).

\emph{The second bound.}  Let \(\Delta\geq\mu>0\) and \(q=\mu/\Delta\in
(0,1]\).  For a subset \(\Omega'\subseteq\Omega\) let \(\mu(\Omega')\) and
\(\Delta(\Omega')\) be the quantities \(\mu\) and \(\Delta\) computed for the
family \((B_i)_{i\in\Omega'}\).  Since
\(\bigwedge_{i\in\Omega}\neg B_i\) implies
\(\bigwedge_{i\in\Omega'}\neg B_i\), the first bound applied to that family
gives
\[
  \Pr\Bigl[\bigwedge_{i\in\Omega}\neg B_i\Bigr]
  \leq\exp\bigl(-\mu(\Omega')+\Delta(\Omega')/2\bigr)
  \qquad\text{for every }\Omega'\subseteq\Omega .
\]
Weight each subset \(\Omega'\) by \(q^{|\Omega'|}(1-q)^{N-|\Omega'|}\);
these weights sum to \(1\).  Each \(i\) lies in subsets of total weight
\(q\), and each ordered pair \(i\neq j\) in subsets of total weight \(q^2\).
So the weighted average of \(\mu(\Omega')-\Delta(\Omega')/2\) is
\[
  q\mu-\frac{q^2\Delta}2=\frac{\mu^2}\Delta-\frac{\mu^2}{2\Delta}
  =\frac{\mu^2}{2\Delta}.
\]
Some \(\Omega'\) has \(\mu(\Omega')-\Delta(\Omega')/2\) at least this
average, and for that \(\Omega'\) the display gives the second bound.
\end{proof}
\end{annexitem}

\begin{annexitem}{annex:p03-w03-p03-065-safe-criterion}{thm:p03-w03-p03-065-safe-criterion}{Almost sure extension statements}
\begin{proof}
Let \((R,H)\) have parameters \(r,v,e\).  Identify the roots with positions
\(1,\ldots,r\) and the nonroots with positions \(1,\ldots,v\).  For a set
\(A\) of nonroot positions, let \(e_A\) be the number of edges of
\(H|_{R\cup A}\) not joining two roots.  For distinct vertices
\(x=(x_1,\ldots,x_r)\) and a tuple \(y=(y_1,\ldots,y_v)\) of distinct
vertices outside \(\{x_1,\ldots,x_r\}\), map root \(i\) to \(x_i\) and
nonroot \(j\) to \(y_j\).  Let \(E_y\) be the set of images of the edges of
\(H\) not joining two roots.  This map is injective, so \(E_y\) consists of
exactly \(e\) distinct pairs.  The tuple \(y\) is an \((R,H)\) extension
over \(x\) exactly when every pair in \(E_y\) is an edge, an event \(B_y\) of
probability \(p^e\).  Assume \(n\geq2(r+v)\) and \(n\geq2\), so
\(p=n^{-\alpha}<1\).

\emph{Not safe implies almost never.}  Suppose \((R,H)\) is not safe.  Then
there is a nonempty set \(A\) of nonroots, \(|A|=a\), such that the
subextension \((R,H|_{R\cup A})\) is not sparse.  Since \(\alpha\) is
irrational, it is dense: \(a-e_A\alpha<0\).  Let \(Z\) be the number of
\((R,H|_{R\cup A})\) extensions over the fixed tuple \((1,\ldots,r)\).  There
are at most \(n^a\) candidate tuples, and each is an extension with
probability \(p^{e_A}\), by the injectivity argument above applied to
\(H|_{R\cup A}\).  So \(\mathbb E[Z]\leq n^ap^{e_A}=n^{a-e_A\alpha}\to0\).
If \(\operatorname{Ext}(R,H)\) holds, the tuple \((1,\ldots,r)\) has an
\((R,H)\) extension, and its entries at the positions in \(A\) form an
\((R,H|_{R\cup A})\) extension over \((1,\ldots,r)\), so \(Z>0\).  By
\cref{thm:p03-w03-p03-065-moments},
\(\Pr[G(n,p)\entails\operatorname{Ext}(R,H)]\leq\Pr[Z>0]\to0\).  For
\(r=0\) the tuple is empty and the argument is unchanged.

\emph{Safe implies almost surely.}  Suppose \((R,H)\) is safe.  If \(e=0\),
then \(\operatorname{Ext}(R,H)\) only asks for \(v\) distinct vertices
outside the \(x\)'s, and holds whenever \(n\geq r+v\).  Assume \(e\geq1\).
For each nonempty set \(A\) of nonroots, \(R\cup A\) is a set \(S\) with
\(R\subsetneq S\subseteq V(H)\), and the subextension on \(S\) has
parameters \(|A|\) and \(e_A\).  By safety \(|A|-e_A\alpha>0\).  Let
\[
  \varepsilon=\min\{|A|-e_A\alpha:\ \emptyset\neq A\subseteq\{1,\ldots,v\}\}>0,
\]
a minimum over finitely many positive numbers.  The case of all \(v\)
nonroots gives \(v-e\alpha\geq\varepsilon\).

Fix distinct \(x_1,\ldots,x_r\), let \(\Omega\) be the set of admissible
tuples \(y\), and let \(N=\sum_{y\in\Omega}\mathbf 1_{B_y}\).  We bound
\(\Pr[N=0]\) by \cref{thm:p03-w03-p03-065-janson}, whose coordinates are
the indicators of the \(\binom n2\) pairs; each has probability \(p<1\), and
\(B_y\) is the event that all coordinates in \(E_y\) equal \(1\).

The expectation is \(\mu=(n-r)(n-r-1)\cdots(n-r-v+1)\,p^e\).  Each factor is
at least \(n-r-v+1\geq n/2\), so
\[
  \mu\geq2^{-v}n^vp^e=2^{-v}n^{v-e\alpha}\geq2^{-v}n^{\varepsilon},
  \qquad
  \mu^2\geq4^{-v}n^{2v}p^{2e}.
\]

Next take \(y\neq y'\) in \(\Omega\) with \(E_y\cap E_{y'}\neq\emptyset\),
and let \(A\) be the set of positions \(j\) such that \(y_j\) is one of
\(y'_1,\ldots,y'_v\).  A pair in \(E_y\cap E_{y'}\) is the image of an edge
\(uw\) of \(H\) not joining two roots.  Any endpoint of \(uw\) that is a
nonroot, at position \(j\), is sent to \(y_j\).  Because the pair also lies
in \(E_{y'}\), the vertex \(y_j\) is one of
\(x_1,\ldots,x_r,y'_1,\ldots,y'_v\), and it is not an \(x\), so \(j\in A\).
Hence \(uw\) is an edge of \(H|_{R\cup A}\) not joining two roots.  At least
one endpoint is a nonroot, so \(A\neq\emptyset\), and
\(|E_y\cap E_{y'}|\leq e_A\).  Therefore
\[
  \Pr[B_y\wedge B_{y'}]=p^{|E_y\cup E_{y'}|}
  =p^{2e-|E_y\cap E_{y'}|}\leq p^{2e-e_A}.
\]
For a fixed nonempty \(A\) with \(|A|=a\), the ordered pairs \((y,y')\) that
lead to \(A\) number at most \(v^an^{2v-a}\).  There are at most \(n^v\)
choices of \(y\).  For \(y'\), there are at most \(v^a\) choices of the
positions in \(y'\) occupied by the vertices \(y_j\), \(j\in A\), and at
most \(n^{v-a}\) choices of the remaining entries.  Hence
\[
  \frac\Delta{\mu^2}
  \leq4^v\sum_{\emptyset\neq A}v^{|A|}n^{-|A|}p^{-e_A}
  =4^v\sum_{\emptyset\neq A}v^{|A|}n^{-(|A|-e_A\alpha)}
  \leq Kn^{-\varepsilon},
  \qquad K=8^vv^v ,
\]
since there are fewer than \(2^v\) sets \(A\) and \(v^{|A|}\leq v^v\).

If \(\Delta<\mu\), the first bound of \cref{thm:p03-w03-p03-065-janson}
gives \(\Pr[N=0]\leq e^{-\mu/2}\leq\exp(-2^{-v-1}n^{\varepsilon})\).  If
\(\Delta\geq\mu\), then \(\mu>0\) and the second bound gives
\(\Pr[N=0]\leq e^{-\mu^2/(2\Delta)}\leq\exp(-n^{\varepsilon}/(2K))\).  With
\(c=\min(2^{-v-1},1/(2K))\), in either case \(\Pr[N=0]\leq
e^{-cn^{\varepsilon}}\).  This bound does not depend on the root tuple.
There are at most \(n^r\) tuples of distinct roots, so by the union bound
\[
  \Pr[G(n,p)\not\entails\operatorname{Ext}(R,H)]\leq n^re^{-cn^{\varepsilon}}
  \to0 .
\]
\end{proof}
\end{annexitem}

\paragraph{Counting extensions.}

In the remaining items \(\alpha\in(0,1)\) is irrational, \(p=n^{-\alpha}\),
and \(n\geq2\), so \(p<1\).  For a rooted graph \((R,F)\) and a tuple \(x\)
of distinct roots, an extension \(y\) over \(x\) is given by the injective
assignment of the positions of \(F\) used in the proof of
\cref{thm:p03-w03-p03-065-safe-criterion}; \(E_y\) is the set of images of
the counted edges, \(|E_y|\) equals the number of counted edges, and \(y\) is
an extension exactly when every pair in \(E_y\) is an edge.

Several proofs bound a sum over sequences \((y^1,\ldots,y^m)\) of tuples over
a fixed \(x\).  The \emph{shape} of such a sequence records which entries of
the \(y^j\) are equal to which.  For fixed \(m\) and fixed numbers of
positions there are finitely many shapes.  The shape determines the number
\(\nu\) of distinct vertices among the entries and the size of
\(U=E_{y^1}\cup\cdots\cup E_{y^m}\); the number of sequences of a given shape
is at most \(n^{\nu}\), and the probability that every \(y^j\) is an
extension is \(p^{|U|}\).  For \(s=1,\ldots,m\) let \(A_s\) be the set of
nonroot positions \(i\) such that \(y^s_i\) is an entry of some earlier
\(y^{s'}\), \(s'<s\).  The entries of \(y^s\) outside \(A_s\) are new
vertices, so \(\nu=\sum_s(w-|A_s|)\) when every \(y^s\) has \(w\) entries.
Every counted edge of \(F\) with an endpoint at a nonroot position outside
\(A_s\) is sent to a pair containing a new vertex, which lies in no earlier
\(E_{y^{s'}}\).  So, writing \(T_s=R\cup A_s\) (positions),
\[
  \nu-\alpha|U|\leq\sum_{s=1}^m\bigl((w-|A_s|)-\alpha\,e_F(T_s)\bigr)
  =\sum_{s=1}^m\lambda_F(T_s).
  \tag{\(\ast\)}
\]

\begin{annexitem}{annex:p03-w03-p03-065-decomposition}{thm:p03-w03-p03-065-decomposition}{Exponent calculus and the rigid--safe decomposition}
\begin{proof}
\emph{Item 1.}  \(V(F)\setminus T\) is the disjoint union of
\(T'\setminus T\) and \(V(F)\setminus T'\).  An edge with an endpoint outside
\(T\) either has an endpoint outside \(T'\), and is counted in \(e_F(T')\),
or has both endpoints in \(T'\) and one outside \(T\), and is counted in
\(e_{F|_{T'}}(T)\); no edge is counted in both.  So
\(e_F(T)=e_{F|_{T'}}(T)+e_F(T')\), and item~1 follows.

\emph{Item 2.}  Let \(R\subseteq T\subsetneq S'\); we show
\(\lambda_{F|_{S'}}(T)<0\).  Put \(T_1=T\cap S\) and \(T_2=T\cup S\).  The
set \(S'\setminus T\) is the disjoint union of \(S\setminus T_1\) and
\(S'\setminus T_2\).  An edge of \(F|_{S'}\) with an endpoint in
\(S'\setminus T_2\) is counted in \(e_{F|_{S'}}(T_2)\); one with both
endpoints in \(S\) and one in \(S\setminus T\) is counted in
\(e_{F|_S}(T_1)\); both kinds are counted in \(e_{F|_{S'}}(T)\), and no edge
is of both kinds.  Hence
\[
  \lambda_{F|_{S'}}(T)\leq\lambda_{F|_S}(T_1)+\lambda_{F|_{S'}}(T_2).
\]
Here \(R\subseteq T_1\subseteq S\) and \(S\subseteq T_2\subseteq S'\).  If
\(T_1\neq S\), the first term is negative by rigidity of \((R,F|_S)\), and the
second is \(0\) if \(T_2=S'\) and negative by rigidity of \((S,F|_{S'})\)
otherwise.  If \(T_1=S\), then \(S\subseteq T\), so \(T_2=T\neq S'\); the
first term is \(0\) and the second is negative.

\emph{Item 3.}  Among the sets \(S\) with \(R\subseteq S\subseteq V(F)\)
such that \(S=R\) or \((R,F|_S)\) is rigid, choose one of largest size; \(R\)
itself is such a set.  Suppose \(S\neq V(F)\) and \((S,F)\) is not safe.
Apply item~3 of \cref{thm:p03-w03-p03-070-weights} at \(\lambda=\alpha\)
in the graph \(F\), with \(X=S\) and \(Y=V(F)\setminus S\); for irrational
\(\alpha\) its notions of rigid and safe over \(X\) are those of
\cref{def:p03-w03-p03-065-dense-sparse} for the rooted graph \((X,F|_{X\cup
Y})\).  It gives a nonempty \(Y_1\subseteq Y\) that is rigid over \(S\); so
\(S'=S\cup Y_1\) satisfies \(S\subsetneq S'\subseteq V(F)\) and
\((S,F|_{S'})\) is rigid.  If
\(S=R\), then \((R,F|_{S'})\) is rigid; otherwise item~2 gives the same.
Either way \(S'\) is a larger set of the kind chosen, a contradiction.
\end{proof}
\end{annexitem}

\begin{annexitem}{annex:p03-w03-p03-065-counting}{thm:p03-w03-p03-065-counting}{Counting theorem for safe extensions}
\begin{proof}
Fix a tuple \(x\) of \(r\) distinct vertices and let \(\Omega\) be the set of
tuples \(y\) of \(v\) distinct vertices outside \(x\), so
\(|\Omega|=(n-r)(n-r-1)\cdots(n-r-v+1)\).  For \(y\in\Omega\) let \(I_y\) be
the indicator that \(y\) is an extension; \(\Pr[I_y=1]=p^e\), so
\(N=N_{R,H}(x)=\sum_yI_y\) has mean \(\mu\).  When \(n\geq2(r+v)\) each factor
of \(|\Omega|\) lies between \(n/2\) and \(n\), which gives
\(2^{-v}n^{v-e\alpha}\leq\mu\leq n^{v-e\alpha}\).  For a set \(A\) of
nonroot positions let \(e_A\) be the number of counted edges of
\(H|_{R\cup A}\), so that \(\lambda_H(R\cup A)=(v-e\alpha)-(|A|-\alpha e_A)\)
by item~1 of \cref{thm:p03-w03-p03-065-decomposition}.  By safety
\[
  \varepsilon_0=\min\{|A|-\alpha e_A:\ A\neq\emptyset\}>0 ,
\]
so \(\lambda_H(R\cup A)\leq v-e\alpha-\varepsilon_0\) for \(A\neq\emptyset\),
including \(A\) the set of all nonroots, where \(\lambda_H=0\).

Fix \(k\geq1\) and expand
\[
  \mathbb E\bigl[(N-\mu)^{2k}\bigr]
  =\sum_{(y^1,\ldots,y^{2k})\in\Omega^{2k}}
  \mathbb E\Bigl[\prod_{j=1}^{2k}(I_{y^j}-p^e)\Bigr].
\]
If some \(E_{y^j}\) is disjoint from every other \(E_{y^i}\), then
\(I_{y^j}-p^e\) depends on coordinates on which the other factors do not
depend, so it is independent of their product and has mean \(0\); the term
vanishes.  Call the other sequences \emph{linked}.  For every sequence,
expanding the product gives \(2^{2k}\) terms
\((-p^e)^{2k-|J|}\prod_{j\in J}I_{y^j}\), \(J\subseteq\{1,\ldots,2k\}\), each
with absolute expectation
\(p^{e(2k-|J|)}p^{|\bigcup_{j\in J}E_{y^j}|}\leq p^{|U|}\); so the term of a
sequence is at most \(2^{2k}p^{|U|}\) in absolute value.

Fix a linked shape.  Join \(i\) and \(j\) when \(E_{y^i}\cap
E_{y^j}\neq\emptyset\).  Every index is joined to another, so each connected
class has at least two indices and there are at most \(k\) classes.  List the
indices class by class, each class in an order in which every index after
the first is joined to an earlier one.  For an index \(s\) that is not first
in its class, a pair of \(E_{y^s}\) lies in an earlier \(E_{y^{s'}}\); the
pair is the image of a counted edge of \(H\), which has a nonroot endpoint at
some position \(q\); the entry \(y^s_q\) lies in that pair, is not a root, so
is an entry of \(y^{s'}\), and \(q\in A_s\).  Thus \(A_s\neq\emptyset\) and
\(\lambda_H(T_s)\leq v-e\alpha-\varepsilon_0\).  For the first index of a
class, \(\lambda_H(T_s)\leq v-e\alpha\), because \(\lambda_H(R\cup A)\leq
v-e\alpha\) for every \(A\).  By \((\ast)\), with \(m=2k\) and \(w=v\),
\[
  \nu-\alpha|U|\leq2k(v-e\alpha)-(2k-k)\varepsilon_0
  =2k(v-e\alpha)-k\varepsilon_0 .
\]
Summing over the finitely many linked shapes, with \(Q_k\) their number,
\[
  \mathbb E\bigl[(N-\mu)^{2k}\bigr]\leq2^{2k}Q_k\,n^{2k(v-e\alpha)-k\varepsilon_0}.
\]
For \(n\geq2(r+v)\), Markov's inequality applied to \((N-\mu)^{2k}\) and
\(\mu\geq2^{-v}n^{v-e\alpha}\) give
\[
  \Pr\bigl[|N-\mu|>\varepsilon\mu\bigr]
  \leq\frac{\mathbb E[(N-\mu)^{2k}]}{(\varepsilon\mu)^{2k}}
  \leq\frac{2^{2k}Q_k4^{kv}}{\varepsilon^{2k}}\,n^{-k\varepsilon_0}.
\]
The bound does not depend on \(x\).  Choose \(k\) with \(k\varepsilon_0>r\);
the bound is then \(o(n^{-r})\).  There are at most \(n^r\) root tuples, so the
probability that some tuple violates \((1-\varepsilon)\mu\leq N\leq
(1+\varepsilon)\mu\) is at most a constant times \(n^{r-k\varepsilon_0}\), which
tends to \(0\).
\end{proof}
\end{annexitem}

\begin{annexitem}{annex:p03-w03-p03-065-rigid-count}{thm:p03-w03-p03-065-rigid-count}{Rigid extensions are boundedly many}
\begin{proof}
Let \(w\) be the number of nonroots of \(F\) and
\(\gamma_0=\min\{-\lambda_F(T):R\subseteq T\subsetneq V(F)\}\), a minimum of
finitely many positive numbers.  Choose an integer \(m_0\geq1\) with
\(m_0\gamma_0>r\), put \(K_0=\binom{wm_0}w\) and \(K=w!\,K_0\).

Fix a tuple \(x\) of distinct roots, and let \(\mathcal Y\) be the family of
\(w\)-element sets of vertices some ordering of which is an extension over
\(x\).  Then \(N_{R,F}(x)\leq w!\,|\mathcal Y|\).  Suppose
\(|\mathcal Y|>K_0\).  Choose members \(Y_1,Y_2,\ldots\) of \(\mathcal Y\) one
at a time, each not contained in the union of those chosen before, for as
long as possible.  If the process stopped after \(s<m_0\) choices, every
member of \(\mathcal Y\) would be a \(w\)-subset of a set of at most \(ws\)
vertices, and \(|\mathcal Y|\leq\binom{ws}w\leq K_0\).  So at least \(m_0\)
choices are made.  Ordering each of \(Y_1,\ldots,Y_{m_0}\) as an extension
gives a sequence \(y^1,\ldots,y^{m_0}\) of extensions over \(x\) in which
each \(y^s\) has an entry not among the entries of earlier ones.

For such a sequence every \(A_s\) misses some nonroot position, so
\(T_s\subsetneq V(F)\) and \(\lambda_F(T_s)\leq-\gamma_0\).  By \((\ast)\)
each such shape contributes at most \(n^{\nu}p^{|U|}\leq n^{-m_0\gamma_0}\) to
the expected number of such sequences, and there are finitely many shapes,
say \(Q\).  By \cref{thm:p03-w03-p03-065-moments},
\(\Pr[|\mathcal Y|>K_0]\leq Qn^{-m_0\gamma_0}\).  Over the at most \(n^r\)
root tuples, the probability that some tuple has more than \(K\) extensions
is at most \(Qn^{r-m_0\gamma_0}\to0\).
\end{proof}
\end{annexitem}

\begin{annexitem}{annex:p03-w03-p03-065-upper-count}{thm:p03-w03-p03-065-upper-count}{An upper bound for every rooted graph}
\begin{proof}
Order the positions of \(S\setminus R\) and of \(V(F)\setminus S\).  Given an
\((R,F)\) extension \(y\) over \(x\), its entries at the positions of
\(S\setminus R\) form an \((R,F|_S)\) extension \(y'\) over \(x\); and its
entries at the positions of \(V(F)\setminus S\) form an \((S,F)\) extension
over the tuple \(xy'\) of \(|S|\) distinct vertices, because every counted
edge of \((S,F)\) is a counted edge of \((R,F)\).  The pair of these two
tuples determines \(y\).  Hence
\[
  N_{R,F}(x)\leq N_{R,F|_S}(x)\cdot\max_{x'}N_{S,F}(x'),
\]
the maximum over tuples \(x'\) of \(|S|\) distinct vertices, where a factor is
read as \(1\) when \(S=R\), respectively \(S=V(F)\).  If \(S\neq R\), then
\((R,F|_S)\) is rigid and, by \cref{thm:p03-w03-p03-065-rigid-count},
almost surely the first factor is at most a constant \(K\) for every \(x\).
If \(S\neq V(F)\), then \((S,F)\) is safe and, by
\cref{thm:p03-w03-p03-065-counting} with \(\varepsilon=1\), almost surely
\(N_{S,F}(x')\leq2\mu_{S,F}\leq2n^{\lambda_F(S)}\) for every \(x'\).  The
intersection of two events of probability tending to \(1\) has probability
tending to \(1\).  So almost surely \(N_{R,F}(x)\leq2Kn^{\lambda_F(S)}\) for
every \(x\), with \(\lambda_F(V(F))=0\) and \(K=1\) when \(S=R\).
\end{proof}
\end{annexitem}

\begin{annexitem}{annex:p03-w03-p03-065-generic-extension}{thm:p03-w03-p03-065-generic-extension}{Generic extension}
\begin{proof}
Let \(P=V(H)\), the roots \(R\) and the nonroot positions
\(1,\ldots,v\), and let \(\varepsilon_0>0\) be as in the proof of
\cref{thm:p03-w03-p03-065-counting}, so that
\(\lambda_{H|_{S}}(R)\geq\varepsilon_0\) whenever \(R\subsetneq S\subseteq
P\), and \(\lambda_H(R)=v-e\alpha\).

\emph{All extensions.}  By \cref{thm:p03-w03-p03-065-counting} with
\(\varepsilon=1/2\), almost surely every tuple \(x\) of distinct roots has
\(N_{R,H}(x)\geq\mu/2\geq2^{-v-1}n^{v-e\alpha}\).

\emph{Extra edges.}  Let \(g\) be a pair of positions of \(P\), not an edge
of \(H\), with at least one nonroot end.  An extension \(y\) over \(x\) has
the corresponding extra edge exactly when \(y\) is an \((R,H+g)\) extension
over \(x\), where \(H+g\) is \(H\) with \(g\) added; and \(y\) violates (2)
exactly when this happens for some \(g\).  Let \(F=H+g\) and let \(S\) be a
set given by item~3 of \cref{thm:p03-w03-p03-065-decomposition} for
\((R,F)\).  Since \(F\) and \(H\) differ only in \(g\),
\(\lambda_F(S)=\lambda_H(S)-\alpha\) if \(g\) has an end outside \(S\), and
\(\lambda_F(S)=\lambda_H(S)\) otherwise.  If \(S=R\), the nonroot end of \(g\)
is outside \(S\), and \(\lambda_F(S)=v-e\alpha-\alpha\).  If \(S\neq R\),
item~1 gives \(\lambda_H(S)=\lambda_H(R)-\lambda_{H|_S}(R)\leq
v-e\alpha-\varepsilon_0\).  In both cases
\(\lambda_F(S)\leq v-e\alpha-\min(\alpha,\varepsilon_0)\).

\emph{Rigid attachments.}  Let \(t\geq1\), and let \(y\) be an extension
over \(x\) that is not \(t\)-generic: there is a set \(Z\) of
\(z\leq t\) vertices, disjoint from the \(x\)'s and \(y\)'s, rigid over
\(X\cup Y\) (the sets of entries of \(x\) and \(y\)), with a vertex adjacent
to some \(y_j\).  List \(Z\) in increasing order as \(z_1,\ldots,z_z\), and
let \(F\) be the graph on \(P\) together with new positions \(1',\ldots,z'\)
whose edges are those of \(H\) and, for every edge of \(G(n,p)\) between a
vertex of \(Z\) and a vertex of \(X\cup Y\cup Z\), the corresponding pair of
positions.  Then \((x,y,z_1,\ldots,z_z)\) is an \((R,F)\) extension over
\(x\), and distinct \(y\) give distinct such extensions.  The parameters of
\((P,F)\) and of its nailextensions count only edges with an end at a
position \(i'\), and these are the images of the edges of \(G(n,p)\) counted
in the rigidity of \(Z\) over \(X\cup Y\); so \((P,F)\) is rigid.  At least
one edge of \(F\) joins some \(i'\) to a nonroot position.  The graphs \(F\)
arising in this way, over all \(z\leq t\), form a finite family
\(\mathcal F_t\), depending only on \(H\) and \(t\), and the number of
extensions over \(x\) violating (3) is at most
\(\sum_{F\in\mathcal F_t}N_{R,F}(x)\).

Fix \(F\in\mathcal F_t\), let \(\gamma_F=\min\{-\lambda_F(T):P\subseteq
T\subsetneq V(F)\}>0\), and let \(S\) be given by item~3 of
\cref{thm:p03-w03-p03-065-decomposition} for \((R,F)\).  Put \(Q=S\cup P\)
and \(S_H=S\cap P\), so \(R\subseteq S_H\subseteq P\).  By item~1,
\(\lambda_F(S)=\lambda_{F|_Q}(S)+\lambda_F(Q)\).  The vertices of \(Q\)
outside \(S\) are those of \(P\setminus S_H\).  The edges of \(F|_Q\) with an
end there include the edges of \(H\) with an end in \(P\setminus S_H\),
numbering \(e_H(S_H)\), and the \(b\) edges of \(F\) between
\(S\setminus P\) and \(P\setminus S_H\).  So
\[
  \lambda_F(S)\leq\lambda_H(S_H)-\alpha b+\lambda_F(Q)
  =(v-e\alpha)-\lambda_{H|_{S_H}}(R)-\alpha b+\lambda_F(Q),
\]
using item~1 for \(H\).  Every term after \(v-e\alpha\) is at most \(0\):
\(\lambda_F(Q)\) is \(0\) if \(Q=V(F)\) and at most \(-\gamma_F\) otherwise,
by rigidity of \((P,F)\).  If \(S_H\neq R\), the second term is at most
\(-\varepsilon_0\).  If \(S_H=R\) and \(Q\neq V(F)\), the last term is at
most \(-\gamma_F\).  If \(S_H=R\) and \(Q=V(F)\), every new position lies in
\(S\) and every nonroot position lies outside \(S\), so \(b\) is the number of
edges of \(F\) between new and nonroot positions, which is at least \(1\).
In every case \(\lambda_F(S)\leq v-e\alpha-\min(\alpha,\varepsilon_0,
\gamma_F)\).

\emph{Conclusion.}  Let \(\gamma\) be the minimum of \(\alpha\),
\(\varepsilon_0\) and the \(\gamma_F\), \(F\in\mathcal F_t\) (only the first
two when \(t=0\)).  By \cref{thm:p03-w03-p03-065-upper-count}, applied to
the finitely many rooted graphs \((R,H+g)\) and \((R,F)\),
\(F\in\mathcal F_t\), there is a constant \(C\) such that almost surely, for
every \(x\), at most \(Cn^{v-e\alpha-\gamma}\) extensions over \(x\) violate
(2) or (3).  Together with the first step, almost surely every \(x\) has at
least \(2^{-v-1}n^{v-e\alpha}\) extensions of which at most
\(Cn^{v-e\alpha-\gamma}\) are bad.  Once \(Cn^{-\gamma}<2^{-v-1}\), some
extension over each \(x\) satisfies (1), (2) and (3).
\end{proof}
\end{annexitem}

\begin{annexitem}{annex:p03-w03-p03-065-shelah-spencer}{thm:p03-w03-p03-065-shelah-spencer}{Zero--one law for \(p=n^{-\alpha}\), \(\alpha\) irrational}
\begin{proof}
The sentences of \(T_\alpha\) (\cref{def:p03-w03-p03-070-theory}) are the
graph axioms, which hold in every \(G(n,p)\); the Nonexistence sentences,
each of which holds almost surely by \cref{thm:p03-w03-p03-068-nonexistence};
and the Generic Extension sentences, each of which holds almost surely by
\cref{thm:p03-w03-p03-065-generic-extension}.  So the hypothesis
\((\mathrm{AS}_\alpha)\) of \cref{def:p03-w03-p03-070-as} holds, and
\cref{thm:p03-w03-p03-072-sparse-zero-one} gives that every first-order graph
sentence has limiting probability \(0\) or \(1\).
\end{proof}
\end{annexitem}
